% ======================================================================
% Apêndices
% ======================================================================
\begin{apendicesenv}
\partapendices

\chapter{Comandos do Git}
\label{ap:git}

Referência dos principais comandos do Git, organizados por categoria. A
documentação completa está em \url{https://git-scm.com/docs}.

\begin{longtable}{|p{3cm}|p{4cm}|p{7cm}|}
    \caption{Comandos do Git} \label{tab:gitcommands} \\
    \hline
    \textbf{Comando} & \textbf{Categoria} & \textbf{Explicação Detalhada} \\
    \hline
    \endfirsthead
    \multicolumn{3}{c}%
    {\textbf{Continuação da Tabela: Comandos do Git}} \\
    \hline
    \textbf{Comando} & \textbf{Categoria} & \textbf{Explicação Detalhada} \\
    \hline
    \endhead
    \multicolumn{3}{|r|}{Continua na próxima página} \\
    \hline
    \endfoot
    \endlastfoot

    git init & Configuração/Setup & Transforma o diretório atual em um repositório Git, criando o diretório \texttt{.git}. Pode ser executado com segurança em um diretório existente sem sobrescrever configurações. \\
    \hline
    git config & Configuração/Setup & Lê ou define variáveis de configuração em nível de sistema, global ou local. Essencial para definir a identidade (user.name, user.email) do autor do commit. \\
    \hline
    git clone [url] & Configuração/Setup & Cria uma cópia local de um repositório remoto. Configura automaticamente a referência 'origin' e faz o checkout da branch principal. \\
    \hline
    git add [file] & Snapshotting Básico & Move alterações de um arquivo do Working Tree para o Index (Staging Area), preparando-o para o próximo commit. \\
    \hline
    git status & Snapshotting Básico & Exibe o estado da Working Tree e do Index, listando arquivos modificados, staged ou não rastreados. \\
    \hline
    git diff & Snapshotting Básico & Mostra as diferenças entre o Working Tree e o Index (alterações não staged). \\
    \hline
    git diff --staged & Snapshotting Básico & Mostra as diferenças entre o Index (Staging Area) e o último commit (HEAD). \\
    \hline
    git commit -m "[msg]" & Snapshotting Básico & Salva o conteúdo atualmente no Index como um novo snapshot permanente (commit) na história. \\
    \hline
    git commit --amend & Manipulação Histórico & Altera o commit anterior, seja modificando sua mensagem ou adicionando/removendo arquivos. Isso reescreve o histórico, gerando um novo SHA. \\
    \hline
    git rm [file] & Gerenciamento Arquivo & Remove um arquivo do Working Tree e do Index. O uso de \texttt{--cached} remove apenas do Index, mantendo o arquivo local. \\
    \hline
    git mv [old][new] & Gerenciamento Arquivo & Move ou renomeia um arquivo de forma rastreada pelo Git. \\
    \hline
    git clean & Gerenciamento Arquivo & Remove arquivos não rastreados (untracked files) do Working Tree. \\
    \hline
    git reset --soft [hash] & Manipulação Histórico & Move o ponteiro HEAD para o commit, mas mantém o Index e o Working Tree intactos (alterações permanecem staged). \\
    \hline
    git reset --mixed [hash] & Manipulação Histórico & (Padrão) Move o HEAD para o commit e reseta o Index (desencena arquivos), preservando o Working Tree. \\
    \hline
    git reset --hard [hash] & Manipulação Histórico & Move o HEAD e reseta o Index e o Working Tree, descartando todas as mudanças locais desde o hash. Altamente destrutivo. \\
    \hline
    git branch & Branching/Navegação & Gerenciamento de branches: lista, cria ou deleta branches locais. \\
    \hline
    git checkout & Branching/Navegação & Comando legado multi-uso. Alterna entre branches ou restaura arquivos antigos/commits, podendo resultar em 'detached HEAD'. \\
    \hline
    git switch & Branching/Navegação & Comando moderno focado em alternar branches. Atualiza a Working Tree e o Index. Utilizado para criar novas branches de forma segura. \\
    \hline
    git merge [branch] & Integração/Merge & Integra alterações de uma branch na atual, criando um 'merge commit' se houver divergência. Operação não-destrutiva. \\
    \hline
    git rebase [base] & Integração/Rebase & Move ou reaplica commits para uma nova base, reescrevendo o histórico para mantê-lo linear. Ideal para branches locais e não publicadas. \\
    \hline
    git rebase -i [base] & Integração/Rebase & Modo interativo do rebase, permitindo squash (combinação), edição ou reordenação de commits. \\
    \hline
    git cherry-pick [hash] & Integração/Portabilidade & Aplica as alterações introduzidas por um único commit específico na branch atual, criando um novo commit equivalente. \\
    \hline
    git revert [hash] & Manipulação Histórico & Cria um novo commit que desfaz as alterações introduzidas por um commit anterior. Usado para desfazer mudanças em histórico compartilhado de forma segura. \\
    \hline
    git fetch & Sincronização Remota & Baixa dados (objetos e refs) de um repositório remoto para o repositório local, sem alterar o Working Tree ou Index (operação segura). \\
    \hline
    git pull & Sincronização Remota & Equivalente a git fetch seguido por uma integração (default: merge). Pode alterar o estado local e causar conflitos imediatamente (operação menos segura). \\
    \hline
    git push [remote][branch] & Sincronização Remota & Carrega commits locais para um repositório remoto. Exige uma operação fast-forward, a menos que \texttt{--force} seja utilizado. \\
    \hline
    git push --tags & Sincronização Remota & Envia tags locais para o repositório remoto. \\
    \hline
    git remote & Sincronização Remota & Gerencia os repositórios remotos rastreados (e.g., listar, adicionar, remover). \\
    \hline
    git log & Auditoria/Inspeção & Exibe o histórico de commits. \\
    \hline
    git shortlog & Auditoria/Inspeção & Fornece um resumo conciso do git log, agrupando commits por autor. \\
    \hline
    git show & Auditoria/Inspeção & Exibe informações detalhadas sobre um objeto Git (commit, tag, etc.). \\
    \hline
    git reflog & Auditoria/Recuperação & Registra as atualizações locais no HEAD e em outras referências, agindo como uma rede de segurança para recuperar commits perdidos após resets ou rebase. \\
    \hline
    git tag & Marcação/Utilitários & Cria, lista, deleta ou verifica objetos de tag, usados para marcar pontos estáticos (releases) no histórico. \\
    \hline
    git tag -a [name] & Marcação/Utilitários & Cria uma tag anotada (com metadados e mensagem), preferida para releases públicas. \\
    \hline
    git stash & Utilitários de Contexto & Salva temporariamente o Working Directory e o Index (alterações ainda sem commit) para permitir a troca de contexto. \\
    \hline
    git stash pop & Utilitários de Contexto & Aplica o último stash salvo e o remove da lista de stashes. \\
    \hline
    git stash apply & Utilitários de Contexto & Aplica o último stash salvo, mas o mantém na lista. \\
    \hline
    git submodule & Utilitários Avançados & Inicializa, atualiza ou inspeciona submódulos (repositórios aninhados). \\
    \hline
    git worktree & Utilitários Avançados & Gerencia múltiplas Working Trees (checkouts) do mesmo repositório, permitindo trabalhar em várias branches simultaneamente. \\
    \hline
    gitk & Utilitários Avançados & O navegador de repositório Git (ferramenta GUI). \\
    \hline
    scalar & Utilitários Avançados & Ferramenta projetada para gerenciar repositórios Git de grande escala (Large Git Repositories). \\
    \hline
    git sparse-checkout & Utilitários Avançados & Reduz a Working Tree para um subconjunto de arquivos rastreados, otimizando o desempenho em repositórios massivos. \\
    \hline
\end{longtable}
\chapter{Comandos do GitHub CLI}
\label{ap:gh}

Referência dos principais comandos do GitHub CLI (\texttt{gh}), organizados
por grupo. A lista completa, sempre atualizada, está no manual oficial
\cite{gh_manual}.

\begin{longtable}{|p{2cm}|p{2.3cm}|p{4cm}|p{6cm}|}
    \caption{Comandos do GitHub CLI (gh)} \label{tab:githubcli_commands} \\
    \hline
    \textbf{Comando Base} & \textbf{Subcomando} & \textbf{Explicação Funcional Detalhada} & \textbf{Exemplo de Sintaxe Chave} \\
    \hline
    \endfirsthead
    \multicolumn{4}{c}%
    {\textbf{Continuação da Tabela: Comandos do GitHub CLI (gh)}} \\
    \hline
    \textbf{Comando Base} & \textbf{Subcomando} & \textbf{Explicação Funcional Detalhada} & \textbf{Exemplo de Sintaxe Chave} \\
    \hline
    \endhead
    \multicolumn{4}{|r|}{\textit{Continua na próxima página}} \\
    \hline
    \endfoot
    \endlastfoot
    
    gh alias & set & Cria um alias para um comando gh, permitindo atalhos personalizados para comandos frequentes & \texttt{gh alias set prc "pr create"} \\
    \hline
    gh alias & list & Lista todos os aliases configurados no GitHub CLI & \texttt{gh alias list} \\
    \hline
    gh alias & delete & Remove um alias previamente configurado & \texttt{gh alias delete prc} \\
    \hline
    gh auth & login & Autentica o usuário no GitHub via navegador web ou token & \texttt{gh auth login} \\
    \hline
    gh auth & logout & Remove a autenticação do usuário atual & \texttt{gh auth logout} \\
    \hline
    gh auth & status & Exibe o status de autenticação atual e usuário conectado & \texttt{gh auth status} \\
    \hline
    gh auth & refresh & Renova a autenticação para um host específico & \texttt{gh auth refresh -{-}hostname github.com} \\
    \hline
    gh auth & token & Exibe o token de autenticação atual & \texttt{gh auth token} \\
    \hline
    gh browse & - & Abre o repositório atual no navegador web & \texttt{gh browse} \\
    \hline
    gh browse & -{-}branch & Abre uma branch específica no navegador & \texttt{gh browse -{-}branch feature-branch} \\
    \hline
    gh browse & -{-}commit & Abre um commit específico no navegador & \texttt{gh browse -{-}commit abc123} \\
    \hline
    gh browse & <número> & Abre a issue ou o pull request de número informado no navegador & \texttt{gh browse 42} \\
    \hline
    gh browse & -{-}settings & Abre as configurações do repositório no navegador & \texttt{gh browse -{-}settings} \\
    \hline
    gh browse & -{-}wiki & Abre a wiki do repositório no navegador & \texttt{gh browse -{-}wiki} \\
    \hline
    gh codespace & code & Abre um codespace no Visual Studio Code & \texttt{gh codespace code} \\
    \hline
    gh codespace & cp & Copia arquivos entre o sistema local e um codespace & \texttt{gh codespace cp local.txt remote:./} \\
    \hline
    gh codespace & create & Cria um novo codespace & \texttt{gh codespace create} \\
    \hline
    gh codespace & delete & Remove um codespace específico & \texttt{gh codespace delete -c my-codespace} \\
    \hline
    gh codespace & jupyter & Abre um codespace no JupyterLab & \texttt{gh codespace jupyter} \\
    \hline
    gh codespace & list & Lista todos os codespaces disponíveis & \texttt{gh codespace list} \\
    \hline
    gh codespace & logs & Exibe os logs de um codespace específico & \texttt{gh codespace logs -c my-codespace} \\
    \hline
    gh codespace & ports & Lista as portas encaminhadas de um codespace & \texttt{gh codespace ports} \\
    \hline
    gh codespace & ports forward & Encaminha uma porta do codespace para o local & \texttt{gh codespace ports forward 3000:4000} \\
    \hline
    gh codespace & ports visibility & Define a visibilidade de uma porta & \texttt{gh codespace ports visibility 3000:public} \\
    \hline
    gh codespace & ssh & Conecta-se a um codespace via SSH & \texttt{gh codespace ssh} \\
    \hline
    gh codespace & stop & Para um codespace em execução & \texttt{gh codespace stop -c my-codespace} \\
    \hline
    gh gist & create & Cria um novo gist a partir de arquivos ou entrada padrão & \texttt{gh gist create script.py} \\
    \hline
    gh gist & clone & Clona um gist específico para o sistema local & \texttt{gh gist clone abc123} \\
    \hline
    gh gist & delete & Remove um gist específico & \texttt{gh gist delete abc123} \\
    \hline
    gh gist & edit & Edita um gist existente & \texttt{gh gist edit abc123} \\
    \hline
    gh gist & list & Lista todos os gists do usuário & \texttt{gh gist list} \\
    \hline
    gh gist & view & Visualiza um gist específico no terminal & \texttt{gh gist view abc123} \\
    \hline
    gh issue & create & Cria uma nova issue no repositório & \texttt{gh issue create -{-}title "Bug" -{-}body "Descrição"} \\
    \hline
    gh issue & list & Lista issues do repositório com filtros opcionais & \texttt{gh issue list -{-}state open} \\
    \hline
    gh issue & status & Mostra o status das issues relevantes para o usuário & \texttt{gh issue status} \\
    \hline
    gh issue & close & Fecha uma issue específica & \texttt{gh issue close 42} \\
    \hline
    gh issue & comment & Adiciona um comentário a uma issue & \texttt{gh issue comment 42 -{-}body "Comentário"} \\
    \hline
    gh issue & delete & Remove uma issue específica & \texttt{gh issue delete 42} \\
    \hline
    gh issue & edit & Edita uma issue existente & \texttt{gh issue edit 42 -{-}title "Novo título"} \\
    \hline
    gh issue & lock & Trava os comentários de uma issue & \texttt{gh issue lock 42} \\
    \hline
    gh issue & reopen & Reabre uma issue fechada & \texttt{gh issue reopen 42} \\
    \hline
    gh issue & transfer & Transfere uma issue para outro repositório & \texttt{gh issue transfer 42 owner/repo} \\
    \hline
    gh issue & view & Exibe detalhes de uma issue específica & \texttt{gh issue view 42} \\
    \hline
    gh project & copy & Copia um projeto para outro dono (usuário ou organização) & \texttt{gh project copy 1 -{-}source-owner @me -{-}target-owner novaorg -{-}title "Cópia" -{-}drafts} \\
    \hline
    gh project & create & Cria um novo projeto & \texttt{gh project create -{-}owner @me -{-}title "Meu Projeto"} \\
    \hline
    gh project & delete & Remove um projeto específico & \texttt{gh project delete 1 -{-}owner @me} \\
    \hline
    gh project & edit & Edita as propriedades de um projeto & \texttt{gh project edit 1 -{-}owner @me -{-}title "Novo Título"} \\
    \hline
    gh project & field-create & Cria um campo personalizado no projeto & \texttt{gh project field-create 1 -{-}owner @me -{-}name "Prioridade" -{-}data-type TEXT} \\
    \hline
    gh project & item-add & Adiciona uma issue ou um pull request ao projeto & \texttt{gh project item-add 1 -{-}owner @me -{-}url \url{https://github.com/owner/repo/issues/1}} \\
    \hline
    gh project & list & Lista projetos disponíveis & \texttt{gh project list -{-}owner owner} \\
    \hline
    gh project & view & Visualiza detalhes de um projeto específico & \texttt{gh project view 1 -{-}owner @me} \\
    \hline
    gh pr & checks & Exibe os status checks de um pull request & \texttt{gh pr checks 15} \\
    \hline
    gh pr & close & Fecha um pull request específico & \texttt{gh pr close 15} \\
    \hline
    gh pr & comment & Adiciona um comentário a um pull request & \texttt{gh pr comment 15 -{-}body "Comentário"} \\
    \hline
    gh pr & create & Cria um novo pull request & \texttt{gh pr create -{-}title "Feature" -{-}body "Descrição"} \\
    \hline
    gh pr & diff & Exibe as diferenças introduzidas pelo pull request & \texttt{gh pr diff 15} \\
    \hline
    gh pr & edit & Edita propriedades de um pull request & \texttt{gh pr edit 15 -{-}title "Novo Título"} \\
    \hline
    gh pr & list & Lista pull requests do repositório & \texttt{gh pr list -{-}state open} \\
    \hline
    gh pr & merge & Mescla um pull request & \texttt{gh pr merge 15 -{-}squash} \\
    \hline
    gh pr & ready & Marca um pull request como pronto para revisão & \texttt{gh pr ready 15} \\
    \hline
    gh pr & reopen & Reabre um pull request fechado & \texttt{gh pr reopen 15} \\
    \hline
    gh pr & review & Adiciona uma revisão a um pull request & \texttt{gh pr review 15 -{-}approve} \\
    \hline
    gh pr & status & Mostra o status dos pull requests relevantes & \texttt{gh pr status} \\
    \hline
    gh pr & view & Exibe detalhes de um pull request específico & \texttt{gh pr view 15} \\
    \hline
    gh pr & checkout & Faz checkout da branch de um pull request & \texttt{gh pr checkout 15} \\
    \hline
    gh release & create & Cria um novo release & \texttt{gh release create v1.0.0 -{-}title "Versão 1.0.0"} \\
    \hline
    gh release & delete & Remove um release específico & \texttt{gh release delete v1.0.0} \\
    \hline
    gh release & download & Baixa os assets de um release & \texttt{gh release download v1.0.0} \\
    \hline
    gh release & list & Lista todos os releases do repositório & \texttt{gh release list} \\
    \hline
    gh release & upload & Faz upload de assets para um release & \texttt{gh release upload v1.0.0 arquivo.zip} \\
    \hline
    gh release & view & Exibe detalhes de um release específico & \texttt{gh release view v1.0.0} \\
    \hline
    gh release & edit & Edita propriedades de um release existente & \texttt{gh release edit v1.0.0 -{-}title "Novo Título"} \\
    \hline
    gh repo & archive & Arquiva um repositório & \texttt{gh repo archive owner/repo} \\
    \hline
    gh repo & clone & Clona um repositório para o sistema local & \texttt{gh repo clone owner/repo} \\
    \hline
    gh repo & create & Cria um novo repositório & \texttt{gh repo create meu-repo -{-}public} \\
    \hline
    gh repo & delete & Remove um repositório & \texttt{gh repo delete owner/repo} \\
    \hline
    gh repo & edit & Edita propriedades de um repositório & \texttt{gh repo edit -{-}description "Nova descrição"} \\
    \hline
    gh repo & fork & Cria um fork de um repositório & \texttt{gh repo fork owner/repo} \\
    \hline
    gh repo & list & Lista repositórios do usuário ou organização & \texttt{gh repo list -{-}limit 10} \\
    \hline
    gh repo & rename & Renomeia um repositório & \texttt{gh repo rename novo-nome} \\
    \hline
    gh repo & sync & Sincroniza um fork com seu repositório upstream & \texttt{gh repo sync} \\
    \hline
    gh repo & view & Exibe detalhes de um repositório & \texttt{gh repo view owner/repo} \\
    \hline
    gh repo & deploy-key & Gerencia chaves de deploy do repositório & \texttt{gh repo deploy-key add chave.pub -{-}title "Servidor"} \\
    \hline
    gh run & cancel & Cancela uma execução de workflow & \texttt{gh run cancel 123456789} \\
    \hline
    gh run & delete & Remove execuções de workflow & \texttt{gh run delete 123456789} \\
    \hline
    gh run & download & Baixa artifacts de uma execução & \texttt{gh run download 123456789} \\
    \hline
    gh run & list & Lista execuções de workflows & \texttt{gh run list} \\
    \hline
    gh run & rerun & Reexecuta um workflow falho & \texttt{gh run rerun 123456789} \\
    \hline
    gh run & view & Exibe detalhes de uma execução & \texttt{gh run view 123456789} \\
    \hline
    gh run & watch & Monitora uma execução em tempo real & \texttt{gh run watch 123456789} \\
    \hline
    gh search & code & Busca por código no GitHub & \texttt{gh search code "função javascript"} \\
    \hline
    gh search & commits & Busca por commits & \texttt{gh search commits "fix bug" -{-}author=user} \\
    \hline
    gh search & issues & Busca por issues e pull requests & \texttt{gh search issues "bug label:bug"} \\
    \hline
    gh search & prs & Busca especificamente por pull requests & \texttt{gh search prs "feature state:open"} \\
    \hline
    gh search & repos & Busca por repositórios & \parbox{3cm}{\texttt{gh search repos "topic:machine-learning"}} \\
    \hline
    gh search & users & Busca por usuários & \texttt{gh search users "nome location:Brasil"} \\
    \hline
    gh secret & list & Lista secrets disponíveis & \texttt{gh secret list} \\
    \hline
    gh secret & remove & Remove um secret específico & \texttt{gh secret remove API\textunderscore KEY} \\
    \hline
    gh secret & set & Define ou atualiza um secret & \texttt{gh secret set API\textunderscore KEY -{-}body "valor"} \\
    \hline
    gh ssh-key & add & Adiciona uma chave SSH à conta & \texttt{gh ssh-key add chave.pub -{-}title "Laptop"} \\
    \hline
    gh ssh-key & list & Lista chaves SSH da conta & \texttt{gh ssh-key list} \\
    \hline
    gh ssh-key & delete & Remove uma chave SSH & \texttt{gh ssh-key delete 123} \\
    \hline
    gh workflow & disable & Desabilita um workflow & \texttt{gh workflow disable "CI Tests"} \\
    \hline
    gh workflow & enable & Habilita um workflow & \texttt{gh workflow enable "CI Tests"} \\
    \hline
    gh workflow & list & Lista workflows disponíveis & \texttt{gh workflow list} \\
    \hline
    gh workflow & run & Executa um workflow manualmente & \texttt{gh workflow run "CI Tests"} \\
    \hline
    gh workflow & view & Exibe detalhes de um workflow & \texttt{gh workflow view "CI Tests"} \\
    \hline
\end{longtable}

% ----------------------------------------------------------------------
\chapter{Glossário}
\label{ap:glossario}

\begin{description}
  \item[Branch] Linha de desenvolvimento dentro do histórico; tecnicamente, um
        ponteiro para um commit.
  \item[Clone] Cópia completa de um repositório remoto no seu computador.
  \item[Commit] Registro de um estado do projeto, com autor, data e mensagem.
  \item[Conflito] Situação em que duas alterações na mesma região de um
        arquivo não podem ser combinadas automaticamente.
  \item[Fork] Cópia de um repositório de outra pessoa na sua conta, para você
        propor mudanças por Pull Request.
  \item[GitHub Actions] Plataforma de automação do GitHub, descrita em
        arquivos YAML.
  \item[GitHub Pages] Hospedagem de sites estáticos a partir de um
        repositório.
  \item[.gitignore] Arquivo que lista o que o Git não deve versionar.
  \item[Merge] Integração de duas linhas de histórico.
  \item[Origin] Nome padrão do repositório remoto principal.
  \item[Pull Request] Pedido para integrar uma branch em outra, com revisão.
  \item[README] Arquivo de apresentação do repositório, em Markdown.
  \item[Remoto] Cópia do repositório em outro lugar, como o GitHub.
  \item[Repositório] Pasta cujo histórico é controlado pelo Git.
  \item[Workflow] Processo automatizado do GitHub Actions.
\end{description}

\end{apendicesenv}
